\documentclass[
    chapter=TITLE,
    12pt, 		% Tamanho da fonte
    openright, 	% Capítulos começam em pág. ímpar (insere página vazia caso preciso)
    oneside, 	% Para impressão somente verso. Oposto a impressão em verso e anverso
    a4paper, 	% Tamanho do papel
    english, 	% Idioma adicional para hifenização
    brazil 		% o último idioma é o principal do documento
    ]{abntex2}
    
\usepackage[utf8]{inputenc}
\usepackage{tikz}
\usepackage{wrapfig}
\usetikzlibrary{calc}
\usepackage{lipsum} 	% para geração de dummy text
\usepackage{multirow}

%% Redefine booleans
%\setboolean{ABNTEXuppersubsection}{false}

% Pacotes
%% Sumário
\usepackage{tocloft}

%% Codificação
\usepackage[T1]{fontenc} 	% Codificação de saída

%% Fonte
\usepackage{times} 			% Times New Roman
\renewcommand
	{\familydefault}
	{\sfdefault} 			% Seta fonte default

\usepackage{lastpage} 		% Usado pela Ficha catalográfica
\usepackage{indentfirst} 	% Indenta o primeiro parágrafo de cada seção.
\usepackage{color} 			% Controle das cores
\usepackage{microtype} 		% Para melhorias de justificação
\usepackage{lscape} 		% Permite a criação de conteúdo em modo paisagem
\usepackage{hyperref} 		% Criação de links
\usepackage{etoolbox}

%% Cabeçalho padrao
\makepagestyle
	{abntheadings}
\makeevenhead
	{abntheadings}
	{\ABNTEXfontereduzida\thepage}
	{}
	{\ABNTEXfontereduzida\textit\leftmark}
\makeoddhead
	{abntheadings}
	{}
	{}
	{\ABNTEXfontereduzida\thepage}
\makeheadrule
	{abntheadings}
	{0pt}
	{\normalrulethickness}
%%%%%%%%%%%%%%%%%%%%%%%%%%%%%%%%%%%%%%%%%%%%%%%%%%%%%%%%%%%%%%%%%


%% Cabeçalho do inicio do capitulo
\makepagestyle
	{abntchapfirst}
\makeoddhead
	{abntchapfirst}
	{}
	{}
	{\ABNTEXfontereduzida\thepage}

%% Tabelas
\usepackage{multirow}

%% Sumário
\usepackage{textcase}

%% Imagens
\usepackage{graphicx} % Inclusão de gráficos

%% Tamanho e alinhamento das legendas
\usepackage[
	singlelinecheck=false,
	justification=justified,
	font=footnotesize
	]
	{caption}

%% Citações
\usepackage[hyphenbreaks]{breakurl}
\usepackage[alf,abnt-emphasize=bf]{abntex2cite}

% Estilos
\autor{Nome SegundoNnome Sobrenome} 	%Nome do autor
% Especificidades para a entrada de autor pessoal: – sobrenome com indicativo de parentesco
% Quando o autor é brasileiro, trate o grau de parentesco como parte do sobrenome.
\newcommand{\entradaAutor}{SOBRENOME, N. S} % Sem ponto no final
% Titulo do trabalho
\titulo{Utilização do Latex na elaboração de trabalhos de conclusão de curso} %Título do trabalho
\newcommand{\englishTitle}{Preparation of Monographs with Latex}
\data{\today}			%Ano de publicação
\local{Nova Friburgo}	%Local
%% Membros da banca
\newcommand{\membrobancaA}{}
\newcommand{\membrobancaB}{}
\newcommand{\membrobancaC}{}

\providecommand
	{\imprimirmembrobancaA}
	{}
\providecommand
	{\imprimirmembrobancaAinst}
	{}
\renewcommand
	{\membrobancaA}
	[2]
	[\imprimirinstituicao]
	{\renewcommand
		{\imprimirmembrobancaA}
		{#2}
		\renewcommand
			{\imprimirmembrobancaAinst}
			{#1}
	}

\providecommand
	{\imprimirmembrobancaB}
	{}
\providecommand
	{\imprimirmembrobancaBinst}
	{}
\renewcommand
	{\membrobancaB}
	[2]
	[\imprimirinstituicao]
	{
		\renewcommand
			{\imprimirmembrobancaB}
			{#2}
			\renewcommand
				{\imprimirmembrobancaBinst}
				{#1}
	}

\providecommand
	{\imprimirmembrobancaC}
	{}
\providecommand
	{\imprimirmembrobancaCinst}
	{}
\renewcommand
	{\membrobancaC}
	[2]
	[\imprimirinstituicao]
	{
		\renewcommand
			{\imprimirmembrobancaC}
			{#2}
			\renewcommand
				{\imprimirmembrobancaCinst}
				{#1}
	}

\newcommand
	{\imprimirmeOrientadorinst}
	{\imprimirinstituicao}
\newcommand
	{\imprimirmeCoorientadorinst}
	{\imprimirinstituicao}

%% Natureza do trabalho
\providecommand{\imprimirnaturezatrabalho}{}
\newcommand{\naturezatrabalho}[1]{
    \renewcommand{\imprimirnaturezatrabalho}{#1}
}

%% Redefinir resumo
\renewenvironment{resumo}[1][\resumoname]{%
   \pretextualchapter{#1}
  }{\PRIVATEclearpageifneeded}
  
%% Redefinir dedicatória
\renewenvironment{dedicatoria}[1][]
	{
	     \ifthenelse{\equal{#1}{}}{
	             \PRIVATEbookmarkthis{\dedicatorianame}
	     }{\preamblealchapter{#1}}
	
	     \vspace*{\fill}
	}
	{}

\renewenvironment{epigrafe}[1][]
{
        \ifthenelse{\equal{#1}{}}{
                \PRIVATEbookmarkthis{\epigraphname}
        }{\pretextualchapter{#1}}

        \vspace*{\fill}
}
{

}

%% Comando para inserir sigla
\newcommand{\sigla}[2][]{
    \item[#1] \textit{#2}
}

%% Redefinição da formatação do \chapter
\renewcommand{\ABNTEXchapterfont}{\normalfont\fontseries{b}\selectfont}
\renewcommand{\ABNTEXchapterfontsize}{\normalsize}
\renewcommand{\ABNTEXpartfont}{\fontseries{b}\selectfont\selectfont}
\renewcommand{\ABNTEXpartfontsize}{\normalsize}
\renewcommand{\ABNTEXsectionfont}{\normalfont\selectfont} %\fontseries{b}
\renewcommand{\ABNTEXsectionfontsize}{\normalsize}
\renewcommand{\ABNTEXsubsectionfont}{\normalfont}
\renewcommand{\ABNTEXsubsectionfontsize}{\normalsize}
\renewcommand{\ABNTEXsubsubsectionfont}{\normalfont}
\renewcommand{\ABNTEXsubsubsectionfontsize}{\normalsize}
\renewcommand{\ABNTEXsubsubsubsectionfont}{\normalfont}
\renewcommand{\ABNTEXsubsubsubsectionfontsize}{\normalsize}
%\renewcommand{\ABNTEXsubsubsectionfont}{\normalfont\fontseries{b}\selectfont}
%\renewcommand{\ABNTEXsubsubsectionfontsize}{\normalsize}
%\renewcommand{\ABNTEXsubsubsubsectionfont}{\normalfont\itshape\selectfont}
%\renewcommand{\ABNTEXsubsubsubsectionfontsize}{\normalsize}



%% Redefinição da formatação de Parágrafos 
\setlength{\parindent}{1.3cm}
\setlength{\parskip}{0cm}

%% Sumário
%\renewcommand{\cftchapterdotsepfont}{\normalfont\normalsize}
\renewcommand*{\cftchapterfont}{\normalfont\fontseries{b}} %\fontseries{b}
\renewcommand{\chapnumfont}{\normalfont\normalsize}
\renewcommand*{\cftsectionfont}{\normalfont} %\fontseries{b}
\renewcommand*{\cftsubsectionfont}{\normalfont}
\renewcommand*{\cftsubsubsectionfont}{\normalfont}
\renewcommand*{\cftsubsubsubsectionfont}{\normalfont}
\renewcommand*{\cftparagraphfont}{\normalfont} %\itshape

%% Forca underline na secao terciaria
\newcommand{\tmpsubsection}[1]{}
\let\tmpsubsection=\subsection
\renewcommand{\subsection}[1]{\tmpsubsection{\underline{#1}}}

%% Forca underline na secao secundaroa
\newcommand{\tmpsection}[1]{}
\let\tmpsection=\section
\renewcommand{\section}[1]{\tmpsection{\textbf{#1}}}

%% Tornar as seções secundários com fonte em maiúscula
%\makeatletter
%\let\oldcontentsline\contentsline
%\def\contentsline#1#2{%
%  \expandafter\ifx\csname l@#1\endcsname\l@section
%    \expandafter\@firstoftwo
%  \else
%    \expandafter\@secondoftwo
%  \fi
%  {%
%    \oldcontentsline{#1}{\MakeTextUppercase{#2}}%
%  }{%
%    \oldcontentsline{#1}{#2}%
%  }%
%}
%\makeatother

%% Tornar a entrada "Referências" com fonte em maiúscula
%%\addto\captionsbrazil{
%%    \renewcommand{\bibname}{REFER\^ENCIAS}
%%}

% Legendas
\makeatletter
\patchcmd{\caption@@@make}
  {\ifcaption@star}
    {\ifcaption@star\small}
    {}{}
\makeatother


\graphicspath{ {figures/} }


\newcommand{\imprimircatalogacao}{%
	\noindent
	UNIVERSIDADE DO ESTADO DO RIO DE JANEIRO\\
	INSTITUTO POLITÉCNICO\\
	CURSO DE ENGENHARIA DE COMPUTAÇÃO\\\\
	
	\noindent
	Reitor: Ruy Garcia Marques\\
	Vice-Reitor: Maria Georgina Muniz Washington\\
	Diretor do Instituto Politécnico: Ricardo Carvalho de Barros\\
	Coordenador de Curso: José Humberto Zani\\
	
	\noindent
	Banca Avaliadora Composta por: \imprimirorientador~(Orientador)\\
	Prof. Dr. Professor 1\\
	Prof. Dr. Professor 2
	
	\begin{center}
		CATALOGAÇÃO NA FONTE \\ UERJ/REDE SIRIUS/BIBLIOTECA CTC/E
		\begin{tikzpicture}[remember picture, overlay]
			\draw[line width = 1pt] ($(current page.north west) + (55mm, -111.5mm)$) rectangle ($(current page.south east) + (-35.5mm,80mm)$);
		\end{tikzpicture}
	\end{center}
	
	Endereço: UERJ - IPRJ \\
			  Rua Bonfim, 25 - Prédio 5, Vila Amélia.\\
			  CEP 28625-570 - Nova Friburgo - RJ - Brasil.\\

	\noindent
	Este trabalho nos termos da legislação que resguarda os direitos autorais é considerado de propriedade da Universidade do Estado do Rio de Janeiro (UERJ). É permitida a transcrição parcial de partes do trabalho, ou mencioná-lo, para comentários e citações, desde que sem propósitos comerciais e que seja feita a referência bibliográfica completa.
	
	\begin{flushright}
		\assinatura{\imprimirautor}
	\end{flushright}
}

\newcommand{\imprimirfolhaaprovacao}{%
	\begin{folhadeaprovacao}
		\begin{center}
			{\large{\imprimirautor}}\\[1cm]
			\vspace{2cm}
			{\Large\bfseries\imprimirtitulo}
		\end{center}		
		\vspace{1cm}
		\hspace{.45\textwidth}
		\begin{minipage}{.5\textwidth}
			\imprimirpreambulo
		\end{minipage}%
		\\\\\\\\
		\vspace{1cm}
		Aprovado em \imprimirdata.\\
		\vspace{1.5cm}
		Banca examinadora:
		
		\assinatura{\imprimirorientador\\Instituto Politécnico - UERJ}
		\assinatura{Professor da banca 1\\Instituto Politécnico - UERJ}
		\assinatura{Professor da banca 2\\Instituto Politécnico - UERJ}
		\begin{center}
			\vfill
			{\large\imprimirlocal}
			\par
			{\large\the\year}
		\end{center}
		
	\end{folhadeaprovacao}
}

\renewcommand{\imprimircapa}{%
	\begin{tikzpicture}[remember picture, overlay]
		\draw[line width = 4pt] ($(current page.north west) + (20mm, -20mm)$) rectangle ($(current page.south east) + (-10mm,10mm)$);
		\draw[line width = 1pt] ($(current page.north west) + (21mm,-21mm)$) rectangle ($(current page.south east) + (-11.5mm,11.5mm)$);
	\end{tikzpicture}
	\begin{capa}
		\begin{minipage}[r]{0.1\linewidth}
			\includegraphics[width=1.5\linewidth]{uerj} 
		\end{minipage}%
		\begin{minipage}[c]{0.7\linewidth}
			\begin{center}
				\textbf{UNIVERSIDADE DO ESTADO DO \\ RIO DE JANEIRO}\\
				\noindent\rule{7cm}{0.4pt}\\
				\textbf{INSTITUTO POLITÉCNICO \\ GRADUAÇÃO EM ENGENHARIA \\DE COMPUTAÇÃO}\\
			\end{center}
		\end{minipage}%
		\begin{minipage}[c]{0.1\linewidth}
			\includegraphics[width=1.5\linewidth]{iprj} 
		\end{minipage}
		
		\vspace*{4.5cm}
		\center
		\textbf{\large\imprimirautor}
		\vspace*{3cm}
		\begin{center}
			\bfseries\LARGE\imprimirtitulo
		\end{center}
		\vfill
		\textbf{
			\large\imprimirlocal
			\\	
			\large \the\year
		}
		\vspace*{1cm}
	\end{capa}
}


\makeatletter
\renewcommand{\folhaderostocontent}{
	\begin{tikzpicture}[remember picture, overlay]
	\draw[line width = 4pt] ($(current page.north west) + (20mm, -20mm)$) rectangle ($(current page.south east) + (-10mm,10mm)$);
	\draw[line width = 1pt] ($(current page.north west) + (21mm,-21mm)$) rectangle ($(current page.south east) + (-11.5mm,11.5mm)$);
	\end{tikzpicture}
	
	\begin{minipage}[r]{0.1\linewidth}
		\includegraphics[width=1.5\linewidth]{uerj} 
	\end{minipage}%
	\begin{minipage}[c]{0.7\linewidth}
		\begin{center}
			\textbf{UNIVERSIDADE DO ESTADO DO \\ RIO DE JANEIRO}\\
			\noindent\rule{7cm}{0.4pt}\\
			\textbf{INSTITUTO POLITÉCNICO \\ GRADUAÇÃO EM ENGENHARIA \\DE COMPUTAÇÃO}\\
		\end{center}
	\end{minipage}%
	\begin{minipage}[c]{0.1\linewidth}
		\includegraphics[width=1.5\linewidth]{iprj} 
	\end{minipage}
	\vspace*{2cm}
	\begin{center}
		{\large\imprimirautor}
		\vspace*{\fill}\vspace*{\fill}
		\begin{center}
			\bfseries\Large\imprimirtitulo
		\end{center}
		\vspace*{\fill}
		\abntex@ifnotempty{\imprimirpreambulo}{%
			\hspace{.45\textwidth}
			\begin{minipage}{.5\textwidth}
				\SingleSpacing
				\imprimirpreambulo
			\end{minipage}%
			\vspace*{2cm}
			\vspace*{\fill}
		}%
		{\abntex@ifnotempty{\imprimirinstituicao}{\imprimirinstituicao
				\vspace*{\fill}}}
		{\large\imprimirorientadorRotulo~\imprimirorientador\par}
		\abntex@ifnotempty{\imprimircoorientador}{%
			{\large\imprimircoorientadorRotulo~\imprimircoorientador}%
		}%
		\vspace*{\fill}
		{\large\imprimirlocal}
		\par
		{\large\the\year}
		\vspace*{1cm}
	\end{center}
}
\makeatother

\preambulo{
	Trabalho de conclusão de curso apresentado como pré-requisito para obtenção do título de Engenheiro de Computação, ao Departamento de Modelagem Computacional, do Instituto Politécnico, da Universidade do Estado do Rio de Janeiro.
}
\orientador{Nome do Orientador}
\tipotrabalho{monografia}

\begin{document}
	%Capa
	\imprimircapa
	
	%Folha de rosto
	\imprimirfolhaderosto
	
	%Catalogação Biblioteca
	\imprimircatalogacao
	
	%Folha de aprovação
	\imprimirfolhaaprovacao

	%% Elemento opcional (Figura 10).
%% A palavra DEDICATÓRIA deve ser grafada em fonte 12, em
%% maiúsculas, negritada e centralizada na parte superior da folha.
%% O texto da dedicatória deve estar localizado na parte inferior da
%% folha, seguindo as regras gerais de apresentação gráfica.

\begin{epigrafe}[Dedicatória]

Dedico este trabalho à minha família, grandes responsáveis pela minha educação.

\end{epigrafe}
	%% A palavra AGRADECIMENTOS deve ser grafada em fonte 12,
%% em maiúsculas, negritada e centralizada na parte superior da folha.
%% O texto dos agradecimentos deve ser separado do título por duas
%% linhas em branco com espaçamento 1,5 e digitado de acordo as regras
%% gerais de apresentação gráfica.
%% Se houver necessidade, o texto pode continuar nas folhas seguin-
%% tes, sem incluir a palavra Agradecimentos.

\begin{agradecimentos}

\lipsum[1]

\end{agradecimentos}
	%% Elemento opcional.
%% É uma citação sem aspas – em fonte 12, estilo normal, com es-
%% paço 1,5 – seguida da indicação de autoria, grafada em fonte 12 e em
%% itálico.
%% O texto deve estar localizado no terço inferior da folha, com o ali-
%% nhamento livre, necessário à epígrafe.

\begin{epigrafe}

\noindent
Lorem ipsum dolor sit amet, consectetur adipiscing elit. Proin nibh est, facilisis et porttitor a, ultrices id tellus.\\
\hspace*{\fill} \textit{Nome Autor}

\end{epigrafe}
	%% 3.1.9 Resumo em língua portuguesa
%% Elemento obrigatório (Figura 14).
%% Consiste na apresentação sucinta dos pontos relevantes do texto,
%% em um único parágrafo. O resumo deve conter entre 150 e 500 pala-
%% vras e fornecer uma visão rápida e clara dos objetivos, da metodologia,
%% dos resultados e das conclusões do trabalho. Na elaboração do resumo,
%% deve-se usar o verbo na voz ativa, na terceira pessoa do singular.

%% Fonte -> TNR ou Arial, corpo 12.
%% A palavra RESUMO deve aparecer em letras maiúsculas
%% e em negrito.
%% O uso de itálico é permitido em palavras estrangeiras.
%% O uso de letras maiúsculas nas palavras-chave
%% restringe-se ao início da palavra, em nomes próprios
%% e siglas, se for o caso.

%% Alinhamento -> A palavra RESUMO deve estar localizada na margem
%% superior da folha e centralizada, e a referência, alinhada
%% à margem esquerda;
%% O alinhamento é justificado para o texto do resumo,
%% que inicia com parágrafo, e para as palavras-chave.

%% Espaçamento -> A palavra RESUMO deve ser separada da referência por
%% duas linhas em branco de 1,5;
%% Espaço 1 na referência e no resumo e, nas palavras-
%% chave, espaço 1,5.

%% Formato do papel,
%% orientação e margens -> Conforme especificado na seção 1.1.

%% Pontuação -> As palavras-chave devem ser separadas por ponto e
%% terminadas por ponto.

\begin{resumo}

\noindent
PINTO, Marcos Belchior Monteiro Morete. \textit{Aplicação de técnicas de Visão Computacional, Machine Learning e SLAM a uma plataforma robótica para realização de tarefas de horticultura}. 2015. \pageref{LastPage} f. Trabalho de Conclusão de Curso (Graduação em Engenharia de Computação) - Instituto Politécnico, Universidade do Estado do Rio de Janeiro, Nova Friburgo, 2015.
\vspace{\onelineskip}

\setlength{\parindent}{1.3cm}
Este trabalho apresenta uma aplicação de métodos de visão computacional e Machine Learning em sistemas embutidos. A visão computacional assume atualmente um importante papel na execução de tarefas repetitivas de inspeção. O Machine Learning vêm sendo aplicado para reconhecimento de padrões em grupo de dados de natureza similar. A aplicação dos sistemas que utilizam visão computacional conjugado com o Machine Learning fundamenta-se  em diferentes áreas da computação moderna. No desenvolvimento do trabalho produziu-se um protótipo para automatização de processos de horticultura de pequeno porte utilizando a plataforma de prototipagem Arduino e o micro computador baseado em um system on a chip (SoC) Raspberry Pi. A combinação das técnicas computacionais e o desenvolvimento de um protótipo favoreceu na comprovação do funcionamento dos algoritmos empregados.		

\vspace{\onelineskip}
\noindent Palavras-chave: Visão computacional. Machine learning. Detecção de objetos. Viola Jones. Classificador de imagens. SLAM.

\end{resumo}
	%% Elemento obrigatório (Figura 15).
%% Consiste em uma tradução do resumo em português para uma
%% língua estrangeira (em inglês, ABSTRACT; em espanhol, RESUMEN;
%% em francês, RÉSUMÉ), em um único parágrafo, seguido das palavras-
%% -chave representativas do conteúdo do trabalho, na língua estrangeira
%% escolhida.
%% O resumo em outra língua também é precedido pela referência
%% do trabalho, substituindo-se o título em português pelo título na língua
%% estrangeira adotada.
%% No caso de teses, é possível incluir dois resumos em língua es-
%% trangeira.
%% A apresentação gráfica e a ordem dos elementos seguem a mes-
%% ma orientação do resumo em português.

\begin{resumo}[Abstract]
\begin{otherlanguage*}{english}

\noindent
\entradaAutor{}. \textit{\englishTitle{}}. 2015. \pageref{LastPage} f. Trabalho de Conclusão de Curso (Graduação em Engenharia de Computação) - Instituto Politécnico, Universidade do Estado do Rio de Janeiro, Nova Friburgo, 2015.
\vspace{\onelineskip}

\setlength{\parindent}{1.3cm}
\lipsum[1]

\vspace{\onelineskip}
\noindent Keywords: Visão computacional. Machine learning. Detecção de objetos. Classificador de imagens.

\end{otherlanguage*}
\end{resumo}
	\listoffigures*
	\cleardoublepage
	\listoftables*
	\cleardoublepage
	%\include{abreviaturas}
	\begin{siglas}
\item[UERJ] Universidade do Estado do Rio de Janeiro
\item[IPRJ] Instituto Politécnico do Rio de Janeiro
\item[HTTP] Hypertext Transfer Protocol
\end{siglas}
	\tableofcontents*
	\cleardoublepage
	
	 % Elementos textuais
	\textual
	\chapter{Cap. 1}

\lipsum[3]

\section{Seção secundária}

\lipsum[1]

\subsection{Seção terciária}

\lipsum[1]

\subsubsection{Seção quartenária}

\lipsum[1]

\subsubsubsection{Seção quinária}

\lipsum[1]


	\chapter{Cap. 2}

\lipsum

\begin{figure}[h]
\centering
\caption{Logo da Unisul}
\includegraphics[width=90mm]{logo_unisul}
\caption*{Fonte: http://unisul.br/}
\end{figure}

\lipsum[1]
	\chapter{Cap. 3}

\lipsum[1]

\begin{table}[h]
\centering
\caption{Taxa de desemprego aberto, por região metropolitana, do primeiro semestre de 1992.}
\begin{tabular}{|c|c|c|c|c|}
\hline
\multirow{2}{*}{Ano e mês} & \multicolumn{4}{|c|}{Região metropolitana}          \\ \cline{2-5} 
                           & Recife & Salvador & Belo Horizonte & Rio de Janeiro \\ \hline
Janeiro                    & 6.10   & 5.43     & 4.77           & 4.24           \\ \hline
Fevereiro                  & 6.44   & 5.18     & 5.00           & 3.81           \\ \hline
Março                      & 6.33   & 5.76     & 5.06           & 4.24           \\ \hline
Abril                      & 6.67   & 6.06     & 4.47           & 4.13           \\ \hline
Maio                       & 6.21   & 7.26     & 4.61           & 5.54           \\ \hline
Junho                      & 5.30   & 6.43     & 4.31           & 3.63           \\ \hline
\end{tabular}
\caption*{Fonte: FUNDAÇÃO INSTITUTO BRASILEIRO DE GEOGRAFIA E ESTATÍSTICA. \textbf{Normas de apresentação tabular}. 3.ed. Rio de Janeiro, 1993. p. 51}
\end{table}

\lipsum[1]
	\chapter{Cap. 4}

Exemplo de citação longa:

\begin{citacao}
\lipsum[1] \cite[5.3]{abntex2cite-alf}
\end{citacao}

Exemplo de citação curta:

Segundo \cite{abntex2cite}, "Lorem ipsum dolor sit amet, consectetur adipiscing elit. Suspendisse adipiscing eu libero non ultricies. Donec accumsan, turpis ut malesuada facilisis, felis augue porta quam, et lacinia turpis ligula et risus.".
	\chapter{Cap. 5}

\lipsum

	% Elementos pós-textuais
	\postextual
	%\include{glossario}
	\apendices

\chapter{Primeiro apêndice}
Primeiro apêndice!

\chapter{Segundo apêndice}
Segundo apêndice
	\anexos

\chapter{Primeiro anexo}
Primeiro anexo!

\chapter{Segundo anexo}
Segundo anexo!
	\bibliography{tcc}
\end{document}